\section{Висновки}

У кваліфікаційній роботі проведено дослідження якості
україно-англійського автоматичного перекладу усного мовлення з акцентом
на семантичну точність та міжфразову когерентність. Дослідження виконано
на корпусі політичного дискурсу (12 сегментів, $\approx$45 хвилин) із
використанням каскадної архітектури ASR (Whisper medium) + MT (OPUS-MT).
Додатково проведено порівняльний аналіз із трьома великими мовними
моделями --- GPT-4o, Claude Sonnet 4, Gemini 2.0 Flash.

\textbf{Основні результати:}

\begin{enumerate}
    \item \textbf{Якість розпізнавання мовлення.} Модель Whisper medium
    продемонструвала WER = 8,62\% та CER = 2,10\%, що є кращим за типові
    показники для української мови (10--15\%). Основні помилки пов'язані
    з розпізнаванням власних назв, географічних назв та абревіатур.

    \item \textbf{Якість машинного перекладу (OPUS-MT).} Спеціалізована
    модель OPUS-MT показала BLEU = 23,12, BERTScore-F1 = 0,397 та
    METEOR = 0,452, що є типовим для морфологічно віддалених мовних пар
    та відповідає очікуваним показникам OPUS-MT на бенчмарку FLORES.

    \item \textbf{Перевага LLM-систем.} Великі мовні моделі (GPT-4o,
    Claude Sonnet 4, Gemini 2.0 Flash) продемонстрували значно вищу
    якість перекладу: BLEU = 55--57 (+147\%), BERTScore-F1 = 0,71--0,74
    (+82\%), METEOR = 0,73--0,75 (+65\%). Різниця статистично значуща
    ($p < 0,001$). Ключовою перевагою LLM є здатність коригувати помилки
    ASR у власних назвах завдяки знанням, отриманим під час навчання.

    \item \textbf{Каскадний ефект.} Виявлено статистично значущу від'ємну
    кореляцію між помилками ASR та якістю перекладу: $r = -0,59$ для
    CER--BLEU ($p = 0,042$) та $r = -0,65$ для CER--BERTScore ($p = 0,022$).
    Це підтверджує гіпотезу про накопичення помилок у каскадному pipeline.

    \item \textbf{Типові помилки.} Ручна анотація виявила 114 помилок
    перекладу OPUS-MT, з яких 79,8\% --- спотворення. Найпроблемніші
    категорії: власні назви (46,2\%), географічні назви (16,5\%), військова
    термінологія (13,2\%). Ідентифіковано 7 критичних помилок, включаючи
    фактичну помилку у валюті та географічні спотворення.

    \item \textbf{Когерентність.} Середня оцінка міжфразової когерентності
    становить 3,2 з 5, що вказує на збереження тематичної структури
    при наявності локальних порушень через непослідовність у передачі
    власних назв.
\end{enumerate}

\textbf{Практичне значення.} Запропонована методологія може застосовуватися
для порівняльного аналізу систем автоматичного перекладу мовлення та для
оцінювання якості перекладу у сферах дипломатичної комунікації. Результати
порівняння OPUS-MT з LLM-системами надають практичні рекомендації щодо вибору
MT-компонента каскадного pipeline: для real-time застосувань --- OPUS-MT з
можливим пост-редагуванням, для офлайн-обробки з вищими вимогами до
якості --- LLM-системи.

\textbf{Обмеження дослідження.} Необхідно зазначити обмеження, що впливають
на інтерпретацію результатів:
\begin{itemize}[nosep]
    \item \textit{Розмір корпусу:} 12 сегментів ($\approx$45 хвилин) є
    достатнім для виявлення основних тенденцій, але обмежує статистичну
    потужність кореляційного аналізу (зокрема, не всі кореляції Спірмена
    досягли рівня значущості).
    \item \textit{Анотація помилок:} ручна анотація виконана одним
    анотатором без обчислення міжанотаторної узгодженості
    (inter-annotator agreement), що обмежує надійність якісного аналізу.
\end{itemize}

\textbf{Напрямки подальших досліджень:} дослідження end-to-end
моделей~\cite{seamlessm4t} для мовної пари uk--en; розширення корпусу
для підвищення статистичної потужності; розробка методів постобробки для
корекції власних назв у каскадному pipeline; порівняння з людським
перекладом; залучення кількох анотаторів для підвищення
надійності якісного аналізу.

\clearpage
